\documentsclass[11pt]{article}
\usepackage{parskip}
\usepackage{color}

\title{Butternut squash and mushroom curry}
\author{Nathaniel Zuk}

\begin{document}
\maketitle

\textcolor{green}{(10/29/2017)}

Based on the recipe in the New York Times: https://cooking.nytimes.com/recipes/1018985-winter-squash-and-wild-mushroom-curry

\textbf{Ingredients:}

\begin{itemize}
\item Butternut squash = 1395 g
\item Mushrooms = 440 g \textcolor{red}{(used mini Bella mushrooms)}
\item Jalapenos = 103 g
\item Scallions = 66 g \textcolor{red}{(in the recipe it said shallots, not scallions)}
\item Cilantro = 32 g
\item Garlic = 17 g
\item Chicken stock ~ 1 cup
\item Coconut milk = 14 oz
\item Tumeric ~ 1 tsp
\item Chili powder ~ 2 tsp \textcolor{red}{(possibly was 1 Tbsp actually...)}
\item Coriander (whole) ~ 1 tsp
\item Cumin ~ 2 tsp
\item Butter ~ 4 Tbsp
\item Salt ~ 2 tsp
\item some ground pepper
\end{itemize}

\textcolor{blue}{I fried the squash, then the mushrooms, jalapenos, scallions, garlic and spices in the large pot for room.  I think this probably also left some of the flavors around in the final dish.  I also ended up using a lot of butter; the original recipe called for oil, but Indian dishes (which this is inspired by) typically use butter. I added the chicken stock to add some more liquid after adding the coconut milk.  After adding the liquids, it took another 10-15 min to finish cooking the squash all the way through, longer than it said in the recipe.}

\textcolor{blue}{The dish came out (surprisingly) good.  I added a decent amount of salt in liberal pinches, but I think it still needed a bit more salt.  I also overdid it on the chili powder, it tastes fine but is a bit too spicy.  I could probably reduce that to 1 tsp in the future.}

\textbf{Servings:}

\begin{itemize}
\item 449 g (10/29)
\item 544 g (10/30)
\item 467 g (10/30)
\item 541 g (10/31)
\item 421 g (10/31)
\item 541 g (11/1)
\item 640 g (11/1)
\end{itemize}

\end{document}